\chapter{Introducción}

El trabajo presentado en esta memoria busca evaluar la utilidad de las herramientas \textit{HDL Coder} y \textit{HDL Verifier} de MATLAB para generar y verificar código HDL a partir de descripciones de alto nivel, con énfasis en la implementación en \textit{FPGAs} y la validación mediante \textit{FPGA-in-the-Loop} (FIL). El foco se sitúa en documentar un flujo de trabajo reproducible que parta desde modelos en \textit{MATLAB} y \textit{Simulink}, continúe con la generación automática de código y concluya con la verificación funcional y temporal sobre hardware objetivo. Se busca complementar la documentación oficial de MATLAB con directrices prácticas y ejemplos que reflejen restricciones de memoria, temporización y arquitectura propias de sistemas embebidos y plataformas reconfigurables.

Se espera aportar evidencia experimental que permita a investigadores e ingenieros evaluar la validez del flujo de diseño propuesto y reducir la brecha entre la descripción de alto nivel y la implementación verificable en hardware. Esto habilita prototipado rápido con garantías funcionales y facilita la simulación con sistemas continuos más cercanos a escenarios reales. Todo el material generado (scripts, modelos, bitstreams y reportes) está organizado para que pueda reutilizarse o extenderse en futuros proyectos.

\section{Motivación y contexto}

Matlab es ampliamente utilizado en la academia y la industria para modelar, simular y analizar algoritmos. Sus entornos de programación (scripts \texttt{.m}, bloques de \textit{Simulink}, \textit{Stateflow}, entre otros) permiten trabajar con distintos niveles de abstracción y acceder a un vasto ecosistema de \textit{toolboxes}. Sin embargo, la verificación funcional en un computador de escritorio no garantiza que el algoritmo pueda ejecutarse eficientemente ni de forma determinista en un sistema embebido o en una FPGA. Las abstracciones de alto nivel facilitan el desarrollo, pero ocultan detalles de recursos, latencia y temporización que resultan críticos cuando se apunta a implementaciones físicas.

Los lenguajes HDL como VHDL o Verilog exigen, en su escritura manual, considerar estructuras de hardware, sincronización, consumo de LUT, registros y bloques DSP, además de la gestión de \textit{clock domains} y \textit{reset}. \textit{HDL Coder} promete traducir modelos de alto nivel a descripciones HDL sintetizables, mientras que \textit{HDL Verifier} habilita la validación contra el modelo original, incluyendo co-simulación con simuladores de terceros y uso de \textit{FIL}. Esta última técnica permite ejecutar el algoritmo directamente en FPGA y mantener el banco de pruebas en MATLAB/Simulink, lo que reduce el tiempo de depuración y acerca la validación a condiciones reales. Además, abre la puerta a probar bloques HDL existentes (de terceros) conectándolos al entorno de referencia para contrastar funcionalidad y temporización sin reescribir el banco de pruebas.

El desafío práctico radica en que no todo modelo de alto nivel es directamente portable a hardware: tipos de datos, dimensiones dinámicas, estructuras de control, acceso a memoria y temporización deben adecuarse a las restricciones del dispositivo. Las capas de abstracción de Matlab no relevan al diseñador de entender estos aspectos, y una configuración inadecuada puede resultar en código HDL que no sintetiza, que no cumple con las restricciones de tiempo o que altera la funcionalidad. Además, el ciclo de iteración hardware--software (generar, sintetizar, cargar, medir) es costoso si no se cuenta con una metodología clara y automatizada.

En el trabajo desarrollado se observó que esta brecha no solo es técnica, sino también metodológica. Muchos errores aparecen por decisiones tomadas en etapas tempranas que parecen menores: un tipo de dato sin saturación explícita, una señal de habilitación omitida o una suposición incorrecta sobre la latencia de salida. Cuando estos detalles no se controlan de forma sistemática, el costo de depuración aumenta y la validación final se vuelve lenta e incierta. Por ello, además de evaluar herramientas, esta memoria enfatiza la necesidad de un proceso ordenado, con trazabilidad de configuraciones y criterios de validación repetibles.

Otro factor relevante es la heterogeneidad de los casos de estudio reales. En un mismo proyecto pueden coexistir bloques de control discreto, módulos puramente secuenciales escritos en HDL y algoritmos de cómputo intensivo con patrones de acceso a memoria no triviales. En ese escenario, la utilidad del flujo no depende únicamente de que funcione en un ejemplo aislado, sino de que permita integrar variantes con exigencias distintas de latencia, ancho de datos y frecuencia interna. Esta motivación orienta la selección de casos y el enfoque comparativo usado en los capítulos siguientes.

\section{Planteamiento del problema}

El objetivo central de esta memoria es diseñar, validar y documentar una metodología para llevar algoritmos descritos en MATLAB/Simulink a implementaciones verificadas en FPGA mediante \textit{HDL Coder} y \textit{HDL Verifier}, apoyándose en \textit{FIL} para cerrar el ciclo de verificación sobre hardware. Además, se incorpora la evaluación de bloques HDL preexistentes (desarrollados fuera de MATLAB) y su verificación mediante \textit{FIL} usando Simulink en un entorno más cercano a la realidad. Los objetivos específicos son:

\begin{itemize}
    \item Definir un flujo que parta en la especificación de alto nivel y culmine en un bitstream funcional, destacando configuraciones clave de \textit{HDL Coder} (tipos de dato fijos, gestión de recursos, \textit{pipelining} y \textit{resource sharing}).
    \item Validar la equivalencia funcional entre el modelo original y el código HDL generado mediante \textit{HDL Verifier}, usando co-simulación y \textit{FIL} para detectar divergencias numéricas o de temporización.
    \item Verificar HDL existentes en un entorno de \textit{Simulink} que emula condiciones cercanas a un sistema real, conectando el bloque a modelos del proceso y usando \textit{FIL}/co-simulación para confirmar su funcionalidad frente a estímulos y respuestas realistas.
    \item Documentar recomendaciones, limitaciones y advertencias que no aparecen de forma explícita en la documentación base, enfocadas en facilitar la extensión a nuevos algoritmos y plataformas.
\end{itemize}

Desde esta perspectiva, el problema puede formularse como una pregunta de ingeniería aplicada: ¿qué condiciones de modelado y verificación permiten que el HDL generado conserve funcionalidad, mantenga un costo de recursos razonable y cumpla la temporalidad requerida por el sistema? Responder esta pregunta exige evaluar el flujo completo, no solo la etapa de generación. Un diseño puede producirse sin errores sintácticos y aun así ser inviable por requerir sobre-reloj excesivo, por no alinearse con el muestreo del banco de pruebas o por introducir retardos incompatibles con el lazo de control.

En consecuencia, la memoria no se limita a reportar ``casos exitosos'', sino que también integra escenarios donde fue necesario ajustar arquitectura, tipos de datos o estrategia de validación para alcanzar una implementación viable. Este enfoque permite que las conclusiones tengan valor práctico para proyectos futuros y evita reducir la discusión a resultados nominales de laboratorio.

\section{Metodología general}

El trabajo se estructura en cuatro etapas principales:

\begin{enumerate}
    \item \textbf{Ajuste del modelo de alto nivel}: se preparan los algoritmos en Matlab/Simulink asegurando tipos fijos, dimensiones estáticas y estructuras de control sintetizables. Se aplican lineamientos de \textit{modeling for HDL} (evitar recursión, matrices dinámicas, \textit{while} sin cotas, etc.).
    \item \textbf{Generación de código HDL}: utilizando \textit{HDL Coder}, se producen descripciones en VHDL/Verilog. Se revisan los reportes de recursos y \textit{critical paths}, y se aplican optimizaciones (p. ej., \textit{pipeline} y repartición de recursos) para cumplir restricciones de área y frecuencia.
    \item \textbf{Verificación funcional y temporal}: con \textit{HDL Verifier}, se realiza co-simulación y se implementa \textit{FPGA-in-the-Loop}. Esta fase comprueba equivalencia funcional respecto al modelo de referencia y detecta desviaciones introducidas por la cuantización o por la latencia de hardware.
    \item \textbf{Documentación}: se miden tiempos de ejecución, utilización de recursos y \textit{throughput}; se comparan con la ejecución en software y se crean guías prácticas en un repositorio Git para su reproducción \cite{pidawRepo}.
\end{enumerate}

Cada etapa está acompañada de scripts automatizados que permiten repetir el flujo con distintas variantes de los algoritmos, registrando configuraciones de \textit{HDL Workflow Advisor}, archivos de constraints y bancos de pruebas usados en las pruebas FIL.

Como referencia de estructura metodológica y de organización del material, también se consideró el repositorio de Jared Soto orientado a \textit{Embedded Coder} \cite{jaredEmbeddedRepo}. Aunque ese trabajo está enfocado en generación de código para microprocesadores (C/C++), su forma de documentar el proceso, ordenar evidencias y presentar resultados fue tomada como guía para este documento, adaptándola al contexto de generación y validación de HDL en FPGA.

Además, la metodología incorpora un criterio de iteración incremental. En lugar de aplicar optimizaciones agresivas desde el inicio, cada caso se desarrolla en tres niveles: versión funcional mínima, versión ajustada para sintetizabilidad y versión optimizada para recursos/temporización. Este orden permite aislar con mayor claridad el efecto de cada decisión y simplifica la depuración cuando aparecen diferencias entre el modelo y el HDL.

También se define un conjunto común de métricas para comparar casos heterogéneos: latencia reportada, sobre-reloj requerido, uso de operadores principales (multiplicadores, sumadores, registros, multiplexores y RAM), y comportamiento funcional en co-simulación/FIL. Con este marco, los resultados del capítulo experimental no dependen de percepciones cualitativas, sino de evidencia comparable entre ejemplos.

\section{Alcances y contribuciones}

El trabajo se acota a la exploración de \textit{HDL Coder} y \textit{HDL Verifier} para la generación y verificación de código HDL orientado a FPGAs. Se tomarán funciones y modelos en Matlab/Simulink que resuelven tareas concretas (p. ej., control, procesamiento de señales, optimización) y se adaptarán al flujo para ilustrar patrones de modelamiento, problemas frecuentes y estrategias de mitigación. Además, se incorporan ejemplos que usan HDL ya existente (IP de terceros) para mostrar cómo integrarlos y validarlos con \textit{FIL} junto a los modelos de referencia.

%La validación experimental se realizó en un conjunto acotado de plataformas representativas de distintas gamas (por ejemplo, placas con dispositivos de entrada como Artix, Cyclone o Zynq) y considerando distintos entornos de ejecución (bare-metal, Linux embebido y, cuando corresponde, aceleradores conectados). Se midió equivalencia funcional mediante co-simulación y FIL, y se registraron tiempos de ejecución y utilización de recursos como línea base. No se optimizó cada diseño para máximo desempeño; los resultados deben interpretarse como referencias iniciales susceptibles de mejora con ajustes específicos de arquitectura.

Las principales contribuciones son:

\begin{itemize}
    \item Evidencia empírica sobre la capacidad de \textit{HDL Coder} para generar código sintetizable y eficiente a partir de modelos de alto nivel, destacando sus límites prácticos.
    \item Un flujo reproducible de verificación con \textit{HDL Verifier} y \textit{FIL} que detecta divergencias funcionales y de temporización antes de comprometer un diseño a producción.
    \item Un repositorio con ejemplos, bitstreams, scripts de automatización y guías breves que facilitan adoptar el flujo en nuevos proyectos, reduciendo tiempos de implementación y riesgo de errores de bajo nivel.
\end{itemize}

Adicionalmente, el trabajo aporta una lectura transversal del flujo, mostrando cómo decisiones de alto nivel se traducen en efectos concretos de implementación. Por ejemplo, se evidencia que la selección de punto fijo no debe entenderse solo como una restricción de herramienta, sino como una decisión de arquitectura que condiciona precisión, latencia y costo de hardware; asimismo, se muestra que el ajuste de \textit{overclocking}/\textit{oversampling} no es un parámetro accesorio, sino una condición necesaria para validar correctamente diseños multicíclo en \textit{FIL}.

En términos de transferencia, la memoria busca dejar una base utilizable por estudiantes e investigadores que necesiten prototipar bloques HDL en plazos acotados, manteniendo un estándar mínimo de rigor experimental. Esto se refleja tanto en la estructura de capítulos como en el material complementario incorporado en anexos.

\section{Organización del informe}

El resto del documento complementa los entregables principales y se organiza como sigue:

\begin{itemize}
    \item El Capítulo 2 resume herramientas de generación automática de código desde alto nivel, con énfasis en \textit{HDL Coder} y \textit{HDL Verifier}.
    \item El Capítulo 3 detalla el flujo propuesto, configuraciones críticas y bancos de prueba usados en co-simulación y \textit{FIL}, junto con las plataformas FPGA evaluadas.
    \item El Capítulo 4 presenta los resultados experimentales: equivalencia funcional, utilización de recursos y tiempos de ejecución comparados.
    \item El Capítulo 5 expone conclusiones y pautas para trabajos futuros, incluyendo oportunidades de optimización y extensiones a otros algoritmos o dispositivos.
\end{itemize}
